\documentclass{a0poster}
\usepackage[utf8]{inputenc}
\usepackage[T1]{fontenc}
\usepackage{geometry}
\usepackage{graphicx}
\usepackage{float}
\usepackage{hyperref}
\usepackage[english,indonesian]{babel}

\title{Optimasi Pemrosesan Template LaTeX Menggunakan Pendekatan Berbasis AI}
\author{Academic Support Agent}
\date{\today}

\begin{document}
\maketitle

\begin{abstract}

\end{abstract}

\chapter{PENDAHULUAN}


\chapter{TINJAUAN PUSTAKA}


% Original content (commented out):
% \documentclass[a0,portrait]{a0poster}
\usepackage{times,colordvi,amsmath,float,color,multicol,caption}
\usepackage{epsfig}
\usepackage{graphics}
\usepackage{graphicx}
\usepackage{scalefnt}
\usepackage[utf8]{inputenc}
\usepackage[brazil]{babel} % ingles / português do Brasil
\usepackage[T1]{fontenc} % fontes com acentos no PDF
\usepackage{amsmath}
\usepackage{amsfonts}
\usepackage{dsfont}
\usepackage{indentfirst}
\usepackage[absolute]{textpos}
\newcommand{\bc}{\begin{center}}
\newcommand{\ec}{\end{center}}
\newcommand{\ex}{\noi {\bf Exemplo} }
\newcommand{\exer}{\noi {\bf Exercício} }
\newtheorem{Def}{Definição}
\newtheorem{pro}{Proposição}
\newtheorem{teo}{Teorema}
\newtheorem{coro}{Corolário}

\usepackage{lipsum}

\newenvironment{Figure}
  {\par\medskip\noindent\minipage{\linewidth}}
  {\endminipage\par\medskip}

\def \ndte{NDT\&E}
\def \rl{Lamb}
%%%%%%%%%%%%%

\setlength{\columnsep}{3cm}  %espaçamento entre as colunas
\setlength{\columnseprule}{0.5pt}  %cor da linha q divide as colunas
\setlength{\parindent}{3cm}
\renewcommand{\baselinestretch}{1.0} %espaçamento entre linhas
\setlength{\parskip}{0cm}
%%%%%%%%%%%%%%%%%%%%%%%%%%%%%%%
%  Kommando-definisjoner:     %
%%%%%%%%%%%%%%%%%%%%%%%%%%%%%%%
% Ingen side-nummerering
\pagestyle{empty}

\definecolor{NTNUBlue}{rgb}{0.0470,0,0.5294}
\definecolor{DarkGreen}{rgb}{0.10,0.40,0.10}
\definecolor{DarkBlue}{rgb}{0.1,0.1,0.5}
\makeatletter
\renewcommand{\abstract}{\@startsection
        {abstract}%                          % the name
        {1}%                                % the level
        {0mm}%                              % the indent
        {-0.6\baselineskip}%                % the beforeskip
        {1mm}%                              % the afterskip
        {\Large\color{NTNUBlue}\bfseries}}% % the style
\renewcommand{\section}{\@startsection
        {section}%                          % the name
        {1}%                                % the level
        {0mm}%                              % the indent
        {-0.6\baselineskip}%                % the beforeskip
        {1mm}%                              % the afterskip
        {\Large\color{NTNUBlue}\bfseries}}% % the style
\renewcommand{\subsection}{\@startsection
        {subsection}%                       % the name
        {2}%                                % the level
        {0mm}%                              % the indent
        {-0.6\baselineskip}%                % the beforeskip
        {1mm}%                              % the afterskip
        {\large\color{NTNUBlue}\bfseries}}% % the style
\renewcommand{\subsubsection}{\@startsection
        {subsubsection}%                    % the name
        {3}%                                % the level
        {0mm}%                              % the indent
        {-0.6\baselineskip}%                % the beforeskip
        {1mm}%                              % the afterskip
        {\large\color{NTNUBlue}\bfseries}}% % the style
\makeatother

\begin{document}

%\DeclareGraphicsExtensions{.jpg}
%\DeclareGraphicsExtensions{.eps}
%%%%%%%%%%%%%%%%%%% Cabeçalho %%%%%%%%%%%%%%%%%%%%%%%%

% \begin{textblock}{4}(0,0.3)
\begin{figure}
\centering
\includegraphics[scale=0.75]{img/logo.pdf}
\label{logo}
\end{figure}
% \end{textblock}


%%%%%%%%%%%%%%%%%%%  Título %%%%%%%%%%%%%%%%%%%%%%%%%%%%

\begin{center}
\parbox{0.7\textwidth}{
\begin{center}
\Huge \sf \vspace{0.1cm} { \textbf{TITULO}}
\end{center}
}
\end{center}
%%%%%%%%%%%%%%%%% Autores %%%%%%%%%%%%%%%%%%%%%

\begin{center}
\vspace{0.3cm}
{\huge{NOME1$^{1*}$,  NOME2$^1$\\
\vspace{0.3cm}
$^1$INST\\
\vspace{0.3cm}
$^{*}$EMAIL\\
}}
\end{center}

\vspace{1cm}
\Large
\begin{multicols}{2}

%%%%%%%%%%%%%%%%%%%%%%%%%%%%%%%%%%%%%%%%%%%%%%%%%%%%%%

\section{Introdução}

\lipsum[2-4]

\section{Metodologia}

\lipsum[1]

\begin{Figure}
    \centering
    \includegraphics[width=\linewidth]{img/example.jpeg}
    \captionof{figure}{Legenda}
    \label{fig:my_label}
\end{Figure}

\section{Resultados}

\lipsum[2-4]

\section{Conclusão}

\lipsum[1-2]

\vspace{1cm}\

\noindent

\begin{thebibliography}{99}

\bibitem[1]{Box}
Box, G. E. P. \& Tiao, G. C. (1973), \textit{Bayesian inference in statistical analysis}, Addison-Wesley Publishing Co., Reading, Mass.-London-Don Mills,
Ont. Addison-Wesley Series in Behavioral Science: Quantitative Methods.

\bibitem[2]{Dunson}
Dunson, D. B. \& Herring, A. H. (2005), 'Bayesian latent variable models for mixed discrete outcomes.', \textit{Biostatistics} \textbf{6}(1), 11-25.


\end{thebibliography}

%%%%%%%%%%%%%%%%%%%%%%%%%%%%%%%%%%%%%%%%%%
\begin{center}
{\parbox{10cm}{

{\bf{
{\color{NTNUBlue}{Agradecimentos}}
}}
}}
\end{center}

\noindent Agradecemos ao XXXXX pelo financiamento desta impressão.


%%%%%%%%%%%%%%%%%%%%%%%%%%%%%%%%%%%%

\end{multicols}

\end{document}

%%%%%%%%%%%%%%%%%%%%%%%%%%%%%%%%%%%%%%


\end{document}